\scanpage{143}
\mspage{88}
% archive sheet number: 142

les $i_X$ ont, \unsure{prr.|par.} à $\unsure{F|\mathcal F}$, la propriété de relèvement à la Quillen.

Je commence par l'hyp. que $C^{\unsure{2|*|?}}$ est déjà un foncteur. Le diagramme

\diagram{\begin{tikzcd}[ampersand replacement=\&]
X\arrow[r,dash,"i_X"]\arrow[d,"u"']\&CX\arrow[d,"C(u)"]\\
X'\arrow[r,dash,"i_{X'}"]\&CX',
\end{tikzcd}}

où $i_X,i_X\in W$, montre que $C$ transforme $W_S$ en $W_S$, donc en des iso de $\mathrm{Ho}^{\unsure{\mathrm{lc}|\mathrm{le}}}_W(S)$. Donc on tient le foncteur cherché $\unsure{\sigma|\gamma}$. On a en fait des foncteurs

\[
\deleted{\mathrm{id}\xrightarrow{\sim}\delta\gamma}
\]

\deleted{à l'aide des}

\diagram{\begin{tikzcd}[ampersand replacement=\&,column sep=large]
\underline{\mathrm{Fib}}\,S\arrow[r,hook,"{\overset{\mathrm{incl.}}{D}}"]\&\deleted{\unsure{\mathrm{C}}}\mathcal M_{/S}\arrow[l,bend left=18,"C"]
\end{tikzcd}}

compatibles avec les localiseurs $W_S$, induisant $\delta$ et $\sigma$, et on a

\[
\mathrm{id}_{\mathcal M_{/S}}\longrightarrow\deleted{\unsure{C}}DC,\qquad\text{de }W_S\ \text{\unsure{obj. par obj.|arg. par arg.}},
\]

\deleted{\unsure{car}} on a \deleted{\unsure{iso}}

\[
X\xrightarrow{i_X}DC(X)=C(X)
\]

avec $i_X\in W_S$. De $=$, on a

\[
\mathrm{id}_{\underline{\mathrm{Fib}}\,S}\longrightarrow CD\qquad\text{de }W_S\ \text{\unsure{obj. par obj.|arg. par arg.}},
\]

\unsure{avec} le même

\[
X\xrightarrow{i_X}CX\qquad(\text{cette fois }X\text{ de }\underline{\mathrm{Fib}}\,S),
\]
